\section*{Author preliminaries and communication model (Section 2)}
For a positive integer $n$, we let $[n]$ denote the set $\{1,\ldots,n\}$.
\paragraph{Binomial Coefficients and Projections.} The binomial coefficient $\binom nk$ equals the number of distinct subsets of $[n]$ with exactly $k$ elements. Similarly, $\binom n{\le k}$ denotes the number of distinct subsets of $[n]$ that have at most $k$ elements. Finally, $\binom{[n]}k$ denotes the set of all subsets of $[n]$ with exactly $k$ elements.
If $x\in\{0,1\}^n$ and $S=\{i_1,\ldots,i_k\}\in\binom{[n]}k$, then $x_S\in\{0,1\}^k$ denotes the projection of $x$ to $S$, given by $x_S=x_{i_1}\cdots x_{i_k}$.
\paragraph{Communication complexity.} We assume basic familiarity with communication complexity \cite{ref38}. We write the definitions here for clarity.
\setcounter{theorem}{4}
\begin{definition}[Protocol]\label{def:protocol}
Let $\mathcal A,\mathcal B,\mathcal Z$ be finite sets. A deterministic protocol $\pi$ over $\mathcal A\times\mathcal B\times\mathcal Z$ is a rooted tree:
\begin{itemize}
\item Each node $v$ is associated with a rectangle $\pi^{-1}(v)=A\times B$, with $A\subseteq X$ and $B\subseteq Y$.
\item Each non-leaf node $v$, associated with a rectangle $\pi^{-1}(v)=A\times B$, is labeled by either (a) a partition $A=\mathop{\dot\bigcup}_{c\in P_v} A_c$ of $A$, in which case we say it is Alice's node or (b) a partition $B=\mathop{\dot\bigcup}_{c\in P_v} B_c$ of $B$, in which case we say it is Bob's node.
\item Each leaf node is labeled by an element of the output domain $\mathcal Z$.
\item The rectangle associated with the root is the input domain $\mathcal A\times\mathcal B$.
\item If a non-leaf node $v$ of Alice is associated with rectangle $\pi^{-1}(v)=A\times B$ and (a) labeled by a partition $A=\mathop{\dot\bigcup}_{c\in P_v} A_c$ of $A$, then for each $c\in P_v$ there will be one child $v_c$ of $v$, which will be associated with the rectangle $\pi^{-1}(v_c)=A_c\times B$; similarly for Bob's nodes.
\end{itemize}
We let the leaf complexity of $\pi$, written $L(\pi)$, be the number of leaf nodes of $\pi$. We let the round complexity of $\pi$, written $R(\pi)$, be the height of $\pi$, i.e., the maximum number of edges in any root-to-leaf path of $\pi$.
A root-to-leaf path $v_1\to\cdots\to v_{k+1}$ is said to have communication length $\sum_{i=1}^k\lceil\log|P_{v_i}|\rceil$ (where $|P_{v_i}|$, as defined above, is the number of parts in the partition of $A$ or $B$ associated with node $v_i$). We then let the communication complexity of $\pi$, written $C(\pi)$, be the maximum communication length of any root-to-leaf path of $\pi$.
Given $(a,b)\in\mathcal A\times\mathcal B$, we let $\pi(a,b)$ denote the (unique) leaf $v$ of $\pi$ having $(a,b)\in\pi^{-1}(v)$. For $z\in\mathcal Z$, we may write $\pi(a,b)=z$ to mean that the leaf $\pi(a,b)$ is labeled by $z$.
A randomized protocol over $\mathcal A\times\mathcal B\times\mathcal Z$ is a distribution over deterministic protocols over $\mathcal A\times\mathcal B\times\mathcal Z$. We will use a boldface Greek letter, such as $\boldsymbol\pi$, to denote a protocol sampled from this distribution. We then let $L(\boldsymbol\pi)$ be the maximum $L(\pi)$ over all $\pi$ in the support of $\boldsymbol\pi$, and likewise for $R(\boldsymbol\pi)$ and $C(\boldsymbol\pi)$.
\end{definition}
\begin{definition}
A function $f:\mathcal A\times\mathcal B\to\mathcal Z$ is said to be computed by a deterministic protocol $\pi$ over $\mathcal A\times\mathcal B\times\mathcal Z$ if we have $f(a,b)=\pi(a,b)$ for every $(a,b)\in\mathcal A\times\mathcal B$. Furthermore, if $\varepsilon\in[0,1]$, then $f$ is said to be computed with error $\varepsilon$ by a randomized protocol $\boldsymbol\pi$ if, for every $(a,b)\in\mathcal A\times\mathcal B$, $\Pr[\boldsymbol\pi(a,b)=f(a,b)]>1-\varepsilon$ (the probability is over the choice of $\boldsymbol\pi$).
We may then define:
\begin{itemize}
\item $L_d^A(f)$ is the minimum leaf complexity $L(\pi)$ among all deterministic protocols $\pi$ that compute $f$, and have round complexity $R(\pi)\le d$, and such that the root node of $\pi$ is Alice's. $L_d^B(f)$ is defined likewise, but for protocols where the root node is Bob's.
\item $L_{d,\varepsilon}^A(f)$ is the minimum $L(\boldsymbol\pi)$ among all randomized protocols $\boldsymbol\pi$ that compute $f$ with error $\varepsilon$, and have $R(\boldsymbol\pi)\le d$, and such that the root node of $\boldsymbol\pi$ is (always) Alice's. $L_{d,\varepsilon}^B(f)$ is defined likewise for randomized protocols where the root node is (always) Bob's.
\item $C_d^A,C_d^B,C_{d,\varepsilon}^A,C_{d,\varepsilon}^B$ are defined analogously where the communication complexity $C(\pi)$ replaces the leaf complexity $L(\pi)$.
\end{itemize}
\end{definition}
Since the communication transcript of a given run of protocol determines the leaf, it must follow that:
\begin{proposition} For any protocol $\pi$, $\log L(\pi)\le C(\pi)$.\end{proposition}
\paragraph{Chromatic number.} All our NP-hardness reductions are from the chromatic number problem:
\begin{definition}[Chromatic number] A coloring of an undirected graph $G$, is a partition of the vertices such that no edge has both endpoints in the same part. The chromatic number of a graph $G$, denoted $\chi(G)$, is the smallest number of parts in a coloring of $G$.\end{definition}
The NP-hardness of approximating the chromatic number has been established by a series of results \cite{ref42,ref24,ref20}, culminating in a paper by Zuckerman \cite{ref65}, where the following was proven:
\begin{theorem}[Hardness For Chromatic Number]\label{thm:chromatic}
For every $\epsilon>0$, there is a deterministic polynomial time algorithm that on an input $x\in\{0,1\}^*$ outputs a graph $G$ on $n$ vertices such that
\begin{itemize}
\item if $x$ is a YES instance of SAT, then $\chi(G)\le n^\epsilon$, and
\item if $x$ is a NO instance of SAT, then $\chi(G)\ge n^{1-\epsilon}$.
\end{itemize}
\end{theorem}
\section*{Author three-round cost characterization (Section 3)}
We start by observing that deterministic 3-round communication complexity may be approximated by a very simple combinatorial quantity.
\begin{definition}
Let $A,B,Z$ be finite sets, and let $M$ be an $A\times B$ matrix over $Z$. Let $\mathcal P=\{P_i\}_{i\in[k]}$ be a partition of (the columns) $B$ for some $k\in\mathbb N$. (That is, $\emptyset\ne P_i\subseteq B$, $\bigcup_{i\in[k]}P_i=B$ and $P_i\cap P_j=\emptyset$ for every $i\ne j$.) For a subset $P\subseteq B$ of columns, we denote by $\operatorname{Cost}_M(P)$ the number of distinct rows of $M$ restricted to columns in $P$, i.e.,
\[\operatorname{Cost}_M(P)=|\{x_P\in Z^P\mid x\in Z^B\text{ is a row of }M\}|.\]
We further define $\operatorname{Cost}_M(\mathcal P)$ to be $\sum_{i=1}^k\operatorname{Cost}_M(P_i)$.
\end{definition}
\begin{proposition}\label{prop:cost}
Let $f:A\times B\to Z$ be a function and $M$ be the $A\times B$ communication matrix (with entries in $Z$) that corresponds to $f$. Let $q\in\mathbb N$ be the maximum number of distinct values in a single row of $M$. We then have
\[ L\le L_3^B(f)\le L\cdot q,\qquad\text{where }L=\min_{\mathcal P}\operatorname{Cost}_M(\mathcal P),\]
and where the minimum is taken over all partitions $\mathcal P$ of $B$. Furthermore we have that
\[\log L\le C_3^B(f)\le\log L+\log q+O(1).\]
\end{proposition}
\begin{proof}
It may be easily seen that $L_2^A(f)$ is lower-bounded by the number of distinct rows in the communication matrix of $f$. Because if Alice's partition of the rows includes a part with two different rows, then Bob's ensuing partition of the columns cannot avoid having a non-monochromatic column. It then follows that, if we have a 3-round protocol $\pi$ where Bob speaks first, and $\mathcal P$ is Bob's partition of the columns in the first round, then $\operatorname{Cost}_M(\mathcal P)$ is a lower-bound on smallest number of leaves that $\pi$ needs to use to compute $f$, and thus $L_3^B(f)\ge L$. The lower-bound on $C_3^B(f)$ now follows from Proposition~7.
Conversely, it is also easy to see that $L_2^A(f)$ is upper-bounded by $q$ times the number of distinct rows in the communication matrix of $f$. The protocol that achieves this bound has Alice tell which kind of row she has, and now in each rectangle the rows are all equal, hence Bob can just tell Alice the color of his column in this row, of which there are $q$ possibilities. Thus $L_3^B(f)\le L\cdot q$, and the same protocol shows that $C_3^B(f)\le\log L+\log q+O(1)$.
\end{proof}
